\input{../header}
\hyphenpenalty=5000

\rhead{Your name: \rule{8cm}{0.15mm}}

\begin{document}

\everymath{\displaystyle}

\section{Learning target T2, version 1}

Use the Taylor series of familiar functions to find the Taylor series for a related function.

\hrulefill


\newcommand{\sinc}{\operatorname{sinc}}
Consider the function $f(x) = \sin(x)$, whose Taylor series is 
\[\sin(x) = \sum_{n=0}^\infty \dfrac{(-1)^n x^{2n+1}}{(2n+1)!} = x - \dfrac{x^3}{3!}+\dfrac{x^5}{5!}-\dfrac{x^7}{7!} + \ldots\]

\begin{enumerate}[leftmargin=15pt]
	\item Write down a Taylor series for the function
	\(\sinc(x) := \dfrac{\sin(x)}{x} = \dfrac{1}{x} \cdot \sin(x).\) \\
	(I am told that this function is important in signal processing for some reason.)

	\vfill

	\item If $g(x) = \sinc(x)$, find the Taylor series for an antiderivative $G(x)$ such that $G(0) = 0$.\\
	(This is called the ``sine integral'' $\operatorname{Si}(x)$ and it is also apparently important in signal processing.)

	\vfill

\end{enumerate}

%%%%%%%%%%%%%%%%%%%%%%%%%%%%%%%%%%%%%%%%%%%%%%%%%%%%%%%
\pagebreak
\section{Learning target T2, version 2}

Use the Taylor series of familiar functions to find the Taylor series for a related function.



\hrulefill

Consider the function $f(x) = \sin(x)$, whose Taylor series is 
\[\sin(x) = \sum_{n=0}^\infty \dfrac{(-1)^n x^{2n+1}}{(2n+1)!} = x - \dfrac{x^3}{3!}+\dfrac{x^5}{5!}-\dfrac{x^7}{7!} + \ldots\]

\begin{enumerate}[leftmargin=15pt]
	\item Write down a Taylor series for the function $\sin(x^2)$.

	\vfill

	\item The Fresnel integral $S(x) = \int \sin(x^2)\ dx$ is important in optics, but as you may surmise, there is not an elementary antiderivative of $\sin(x^2)$. Darn! Guess we'll have to write down a power series representation for it instead.

	\vfill

\end{enumerate}
%%%%%%%%%%%%%%%%%%%%%%%%%%%%%%%%%%%%%%%%%%%%%%%%%%%%%%%

\end{document}